\[ %%%%%%%%%%%%%%%%%%%%%%%%%%%%%%%%%%%%%%%%%%%%%%%%%%%%%%%%%%%%%%%%%%%%%%%%%%%%%%%% \subsection{Multi-Step Forecasting Loss} \label{sec:multistep_loss} %%%%%%%%%%%%%%%%%%%%%%%%%%%%%%%%%%%%%%%%%%%%%%%%%%%%%%%%%%%%%%%%%%%%%%%%%%%%%%%% Long-horizon prediction represents one of the most challenging problems in learning nonlinear dynamical systems. Although many neural architectures can achieve accurate short-term forecasting, small prediction errors may accumulate during recursive rollout, resulting in divergence from the true physical trajectory. PHYSGEN addresses this challenge by explicitly optimizing the model for multi-step trajectory prediction rather than only one-step state reconstruction. %%%%%%%%%%%%%%%%%%%%%%%%%%%%%%%%%%%%%%%%%%%%%%%%%%%%%%%%%%%%%%%%%%%%%%%%%%%%%%%% \subsubsection{Recursive Forecasting Formulation} %%%%%%%%%%%%%%%%%%%%%%%%%%%%%%%%%%%%%%%%%%%%%%%%%%%%%%%%%%%%%%%%%%%%%%%%%%%%%%%% Given an initial system state $X_t$, the predicted future trajectory is obtained through recursive application of the learned evolution operator: \begin{equation} \hat{X}_{t+k} = \Phi_{\theta}^{k} ( X_t ), \end{equation} where $k$ represents the forecasting horizon. The predicted trajectory is therefore: \begin{equation} \hat{\mathcal{T}} = \{ \hat{X}_{t+1}, \hat{X}_{t+2}, \dots, \hat{X}_{t+K} \}. \end{equation} The objective is to minimize the deviation between the predicted trajectory and the ground truth physical evolution: \begin{equation} \mathcal{T} = \{ X_{t+1}, X_{t+2}, \dots, X_{t+K} \}. \end{equation} %%%%%%%%%%%%%%%%%%%%%%%%%%%%%%%%%%%%%%%%%%%%%%%%%%%%%%%%%%%%%%%%%%%%%%%%%%%%%%%% \subsubsection{Horizon-Weighted Forecasting Loss} %%%%%%%%%%%%%%%%%%%%%%%%%%%%%%%%%%%%%%%%%%%%%%%%%%%%%%%%%%%%%%%%%%%%%%%%%%%%%%%% The multi-step forecasting loss is defined as: \begin{equation} \mathcal{L}_{forecast} = \sum_{k=1}^{K} w_k \| X_{t+k} - \hat{X}_{t+k} \|^2, \end{equation} where $w_k$ represents the contribution weight of the prediction horizon. The weighting strategy can be selected according to the physical task. For example: \begin{equation} w_k = \frac{k^\gamma} {\sum_{i=1}^{K}i^\gamma}, \end{equation} where $\gamma$ controls the emphasis on long-term prediction. A positive value of $\gamma$ increases the importance of distant future states. %%%%%%%%%%%%%%%%%%%%%%%%%%%%%%%%%%%%%%%%%%%%%%%%%%%%%%%%%%%%%%%%%%%%%%%%%%%%%%%% \subsubsection{Teacher-Forced and Free-Running Prediction} %%%%%%%%%%%%%%%%%%%%%%%%%%%%%%%%%%%%%%%%%%%%%%%%%%%%%%%%%%%%%%%%%%%%%%%%%%%%%%%% During training, many sequence models use teacher forcing: \begin{equation} \hat{X}_{t+k} = \Phi_{\theta} ( X_{t+k-1} ), \end{equation} where the true previous state is provided to the model. Although effective for optimization, teacher forcing creates a discrepancy between training and inference conditions. During deployment, the model must operate autonomously: \begin{equation} \hat{X}_{t+k} = \Phi_{\theta} ( \hat{X}_{t+k-1} ). \end{equation} This difference is known as exposure bias. PHYSGEN reduces this effect by incorporating free-running rollout during training: \begin{equation} \hat{X}_{t+k} = \Phi_{\theta}^{k}(X_t). \end{equation} %%%%%%%%%%%%%%%%%%%%%%%%%%%%%%%%%%%%%%%%%%%%%%%%%%%%%%%%%%%%%%%%%%%%%%%%%%%%%%%% \subsubsection{Trajectory Consistency Loss} %%%%%%%%%%%%%%%%%%%%%%%%%%%%%%%%%%%%%%%%%%%%%%%%%%%%%%%%%%%%%%%%%%%%%%%%%%%%%%%% Accurate long-term prediction requires preserving the structure of the complete trajectory rather than minimizing independent state errors. Therefore, PHYSGEN introduces trajectory consistency: \begin{equation} \mathcal{L}_{traj} = \| \Delta \hat{X}_{t:t+K} - \Delta X_{t:t+K} \|^2, \end{equation} where: \begin{equation} \Delta X_t = X_{t+1}-X_t \end{equation} represents the system transition. This term encourages the model to learn the underlying dynamics rather than only matching individual states. %%%%%%%%%%%%%%%%%%%%%%%%%%%%%%%%%%%%%%%%%%%%%%%%%%%%%%%%%%%%%%%%%%%%%%%%%%%%%%%% \subsubsection{Physics-Guided Rollout Constraint} %%%%%%%%%%%%%%%%%%%%%%%%%%%%%%%%%%%%%%%%%%%%%%%%%%%%%%%%%%%%%%%%%%%%%%%%%%%%%%%% Long-horizon trajectories must remain physically admissible during recursive prediction. Therefore, PHYSGEN introduces a rollout physics constraint: \begin{equation} \mathcal{L}_{rollout}^{phys} = \sum_{k=1}^{K} \| \mathcal{C} ( \hat{X}_{t+k} ) \|^2, \end{equation} where $\mathcal{C}$ represents physical consistency conditions. Examples include: \begin{itemize} \item conservation violations; \item unstable energy growth; \item invalid parameter ranges; \item violation of boundary conditions. \end{itemize} %%%%%%%%%%%%%%%%%%%%%%%%%%%%%%%%%%%%%%%%%%%%%%%%%%%%%%%%%%%%%%%%%%%%%%%%%%%%%%%% \subsubsection{Long-Horizon Optimization Objective} %%%%%%%%%%%%%%%%%%%%%%%%%%%%%%%%%%%%%%%%%%%%%%%%%%%%%%%%%%%%%%%%%%%%%%%%%%%%%%%% The complete forecasting objective combines state accuracy, trajectory consistency, and physical validity: \begin{equation} \boxed{ \mathcal{L}_{multi} = \mathcal{L}_{forecast} + \lambda_t \mathcal{L}_{traj} + \lambda_p \mathcal{L}_{rollout}^{phys} } \end{equation} This formulation directly optimizes the behavior of PHYSGEN during long-term autonomous prediction. %%%%%%%%%%%%%%%%%%%%%%%%%%%%%%%%%%%%%%%%%%%%%%%%%%%%%%%%%%%%%%%%%%%%%%%%%%%%%%%% \subsubsection{Error Growth Analysis} %%%%%%%%%%%%%%%%%%%%%%%%%%%%%%%%%%%%%%%%%%%%%%%%%%%%%%%%%%%%%%%%%%%%%%%%%%%%%%%% For nonlinear dynamical systems, prediction errors may grow according to: \begin{equation} \epsilon_{k} \leq C e^{\lambda k} \epsilon_0, \end{equation} where $\lambda$ represents an effective instability factor. The objective of PHYSGEN is to reduce effective error amplification through: \begin{itemize} \item physics-guided latent evolution; \item adaptive stability control; \item constrained graph propagation. \end{itemize} Therefore, the model aims to achieve: \begin{equation} \lambda_{PHYSGEN} < \lambda_{data-only}. \end{equation} %%%%%%%%%%%%%%%%%%%%%%%%%%%%%%%%%%%%%%%%%%%%%%%%%%%%%%%%%%%%%%%%%%%%%%%%%%%%%%%% \subsubsection{Summary} %%%%%%%%%%%%%%%%%%%%%%%%%%%%%%%%%%%%%%%%%%%%%%%%%%%%%%%%%%%%%%%%%%%%%%%%%%%%%%%% The multi-step forecasting loss establishes long-horizon prediction as the primary optimization target of PHYSGEN. By explicitly training recursive trajectories, incorporating physics-guided rollout constraints, and reducing exposure bias, PHYSGEN is designed to maintain stable and physically consistent predictions beyond short forecasting horizons. \] 
